\section{结果}\label{sec:results}

% 呈现方式（赛题补充说明四）：以图为主；「改之前 vs 改之后」必须同图同轴可比。

\subsection{当前可跑的基线（ExpertPolicy 驱动，非学习后）}
\todo{填入：\texttt{exp\_arena\_expert\_ref\_v2}（3 seed $\times$ 12 鱼 $\times$ 600 步）——}
\todo{生存 0.9472$\pm$0.0459、捕食 0.5407$\pm$0.1180、逃脱 0.2292$\pm$0.0625、}
\todo{能量率 $-1.031\times10^{-3}$、fitness 0.5123$\pm$0.0428。}
\todo{再次强调：以上为「个体比值均值」口径。}


\subsection{已跑通的一组基线}
\begin{figure}[t]
  \centering
  \includegraphics[width=0.95\linewidth]{fig_metrics_mean_std}
  \caption{Arena 主指标（ExpertPolicy 驱动，3 seed \ $\times$ 12 鱼 $\times$ 600 步）。
  左：0--1 比值类；右：\texttt{energy\_efficiency} 为\textbf{每步速率}（$\sim10^{-3}$），
  因量纲不同单列一轴。误差棒为跨 seed 的 mean $\pm$ std。}
  \label{fig:metrics}
\end{figure}

\begin{figure}[t]
  \centering
  \includegraphics[width=0.95\linewidth]{fig_env_absolute}
  \caption{为何比值会掩盖猎物密度效应：猎物变多时\textbf{分子与分母同时上升}
  （food\_rich 捕获 55$\to$100，遭遇 722$\to$2325），故比值反而下降。}
  \label{fig:env-absolute}
\end{figure}

\subsection{环境对照}
\todo{predator\_rich 的 escape 效应量 $d=+2.89$、prey\_capture $d=+1.92$；}
\todo{food\_rich 绝对捕获 55$\to$100 但比值反降（分母 722$\to$2325）。}

\subsection{学习前后对照}
\todo{\textbf{阻塞}：BC \emph{模块与契约}已就绪（\texttt{src/evogenesis/learning/}），
但\textbf{未接入流水线}（\texttt{pipeline/模型链装配.md} \S4「未接」），
且当前评估仍由 \texttt{ExpertPolicy} 驱动 $\Rightarrow$ 「学习前 / 学习后」对照暂无数据。
\textbf{不得用其他数字顶替}。}

\subsection{典型案例}
\todo{补充说明四要求「能直接看见」：选一条成功、一条失败的行为轨迹做可视化。}
